(a) Send the germ of $f$ at $\varphi(P)$ to the germ of $f\circ\varphi$ at $P$. This is well-defined because inverse images of neighborhoods are neighborhoods, and composition preserves regularity. It is a $k$-algebra homomorphism and takes the maximal ideal into the maximal ideal, since $(f\circ\varphi)(P)=f(\varphi(P))$.

(b) An isomorphism clearly has the stated properties. Conversely, let $\psi=\varphi^{-1}$ be the continuous inverse. For a regular function $f$ near $P$, surjectivity of $\varphi_P^*$ gives a regular germ $g$ near $\varphi(P)$ with $f=g\circ\varphi$ near $P$. As $\varphi$ is a homeomorphism, this means $f\circ\psi=g$ on a neighborhood of $\varphi(P)$. Thus $\psi$ is a morphism.

(c) If a germ pulls back to zero on a nonempty open $U\subseteq X$, its representative vanishes on $\varphi(U)$. This image is dense in $Y$: every nonempty open $V\subseteq Y$ has nonempty open inverse image, which meets $U$. Hence the representative vanishes on its domain, so the germ is zero.
